\documentclass[10pt,a4paper]{amsart}
\usepackage[margin=19mm]{geometry}
\usepackage{amsmath,amssymb,mathtools}
\usepackage[scr=rsfs,cal=euler]{mathalfa}
\usepackage{xfrac}
\usepackage[unicode,hidelinks,bookmarks=false]{hyperref}
\newcommand{\ie}{i.e.\ }
\newcommand{\cf}{cf.\ }
\newcommand{\resp}{respectively\ }
\newcommand{\Subsection}{\subsection}
\providecommand{\nobreakdash}{\nobreak\hbox{-}}
\setlength{\emergencystretch}{2em}
\multlinegap0pt
\usepackage[shortlabels]{enumitem}
\usepackage{amssymb}
\usepackage{mathtools}
\usepackage{xcolor}
\newtheorem{lemma}{Lemma}
\newtheorem{theorem}{Theorem}
\newcommand{\NA}{\operatorname{NA}}
\newcommand{\R}{\mathbb R}
\newcommand{\norm}[1]{\left\lVert #1\right\rVert}
\title{Smooth Diffeomorphisms and Mahler's Problem on Liouville Numbers}
\author{Diego Marques}
\date{Source: arXiv:2609.00376v3 (CC BY 4.0)}
\begin{document}
\maketitle

\noindent\emph{Source and licence.} Smooth Diffeomorphisms and Mahler's Problem on Liouville Numbers. Version arXiv:2609.00376v3 (CC BY 4.0), distributed by arXiv under the Creative Commons Attribution 4.0 International licence, http://creativecommons.org/licenses/by/4.0/, as recorded in the arXiv licence metadata for this exact version and shown on the abstract page https://arxiv.org/abs/2609.00376v3. Full licence notice: \texttt{licenses/arxiv-2609.00376-CC-BY-4.0.txt}.\par\smallskip

\noindent\textit{Editorial scope.} Counted unit: the article's theorem thm:sewing (source0.tex lines 640--657). The article's complete original proof of it is delivered with it (source0.tex lines 659--798, one \texttt{proof} environment of 140 lines). No mathematics is written, repaired or re-derived: the statement and the proof are the article's own lines, in the article's own notation, and no retained line is substituted or re-flowed.  Retained uncounted as the source-local context the proof needs: lines 469--492; lines 525--536; lines 596--604.  Nothing was omitted from any retained span, so the package carries no omission record.  No retained line contains a citation, so the wrapper carries no bibliography and no bibliography entry.  No retained line contains a figure environment, an image-inclusion command or drawing code, so no image is delivered and the package declares no asset dependency.  Editorial formatting only, in addition: an AMS \texttt{amsart} preamble that restates the article's own declarations and macros used by the retained lines, namely the environments \texttt{lemma}, \texttt{theorem} on separate counters, together with the standard packages those lines need.  The retained environments print in this document as Lemma 1, Theorem 1 in order of appearance, whereas the article prints the same items with the numbering of its own full document.  The complete native source of the article as distributed in the arXiv e-print for this exact version is bundled unchanged as \texttt{sources/209019/source0.tex}.
\begin{lemma}[Sewing endpoint jets]\label{lem:stitch}
Let $r\geq0$, let $J=[a,b]$ be a nondegenerate compact interval, let $P$ be a polynomial,
and let $\eta>0$. Then one can choose $\delta>0$ such that, whenever real
numbers $\alpha_j,\beta_j$ are prescribed for $0\leq j\leq r+1$ and satisfy
\[
|\alpha_j-P^{(j)}(a)|<\delta,
\qquad
|\beta_j-P^{(j)}(b)|<\delta
\qquad(0\leq j\leq r),
\]
there exists a polynomial $S$ such that
\[
S^{(j)}(a)=\alpha_j,
\qquad
S^{(j)}(b)=\beta_j
\qquad(0\leq j\leq r+1),
\]
and
\[
\norm{S-P}_{C^r(J)}<\eta.
\]
No bound is required on the order-$(r+1)$ data
$\alpha_{r+1},\beta_{r+1}$.
\end{lemma}

\begin{lemma}[Approximation over a dense subfield]\label{lem:dense-field}
Let $F\subset\R$ be a dense subfield, let $P\in\R[x]$, let $J$ be a compact
interval, let $r,N_0\geq0$, and let $\delta>0$. Then, for every integer
\[
N>\max\{\deg P,N_0\},
\]
there exists a polynomial $Q\in F[x]$ of degree exactly $N$ such that
\[
\norm{Q-P}_{C^r(J)}<\delta.
\]
In particular, one may choose $Q$ with $\deg Q\geq N_0$.
\end{lemma}

\begin{lemma}[Local abundance of gaps]\label{lem:gaps}
Let $x\in E$, let $V$ be an open interval containing $x$, and let $N\geq1$.
Then there exists $j\geq N$ such that
\[
\overline{U_j}\subset V.
\]
In fact, $V$ contains the closures of infinitely many complementary
components of $\R\setminus E$.
\end{lemma}

\begin{theorem}[Arithmetic polynomial sewing]\label{thm:sewing}
Let $E\subset\R$ be nonempty, compact, perfect, and nowhere dense, and let
$F\subset\R$ be a dense subfield. Then there exists a function
$g\in C_c^\infty(\R)$ such that
\begin{enumerate}[label=\textup{(\roman*)}]
\item for every connected component $I$ of $\R\setminus E$, there exists
$P_I\in F[x]$ such that
\[
g|_I=P_I|_I;
\]
\item $g$ is real-analytic on $\R\setminus E$;
\item $g$ is not real-analytic at any point of $E$.
\end{enumerate}
Consequently,
\[
\NA(g)=E.
\]
\end{theorem}

\begin{proof}
Write $a=\min E$ and $b=\max E$, and enumerate the bounded gaps as
$U_1,U_2,\dots$. We construct functions $g_n$ and polynomials
$Q_1,\dots,Q_n\in F[x]$ inductively so that
\begin{enumerate}[label=\textup{(A\arabic*)},leftmargin=3.1em]
\item $g_n\in C^n(\R)$ and $g_n=0$ on $\R\setminus[a,b]$;
\item $g_n|_{U_j}=Q_j|_{U_j}$ for $1\leq j\leq n$;
\item for every $J\in\mathcal J_n$, the restriction $g_n|_J$ is a real polynomial;
\item $\deg Q_1<\cdots<\deg Q_n$;
\item
\begin{equation}\label{eq:increment}
\norm{g_{n+1}-g_n}_{C^n([a,b])}\leq2^{-n-2}
\qquad(n\geq0).
\end{equation}
\end{enumerate}
For $n=0$, condition \textup{(A4)} is vacuous, and we set $g_0=0$.

Assume that $g_n$ has been constructed. Write
\[
U_{n+1}=(c,d).
\]
This gap lies in the interior of a unique interval
$J_*\in\mathcal J_n$, and on $J_*$ we have $g_n=P_*$ for some real
polynomial $P_*$. Set
\[
\eta_n:=2^{-n-2}.
\]
We shall choose the polynomial $Q_{n+1}\in F[x]$ to be assigned to the new gap
$U_{n+1}$, close to $P_*$ through order $n$ on $[c,d]$, while requiring its
degree to exceed those of all previously assigned gap polynomials.

Let $K=[u,v]\in\mathcal J_{n+1}$, let $J\in\mathcal J_n$ be its unique
parent, and write $g_n|_J=P_J$. The polynomial to be placed on $K$ must
match, through order $n+1$, the piece lying across each endpoint. That
neighboring piece is one of:
\begin{itemize}
\item the zero polynomial, at $a$ or $b$;
\item a previously assigned polynomial $Q_j$, at an old gap endpoint;
\item the new polynomial $Q_{n+1}$, at $c$ or $d$.
\end{itemize}
Denote the corresponding target jets at $u$ and $v$ by
\[
\alpha_0,\dots,\alpha_{n+1},
\qquad
\beta_0,\dots,\beta_{n+1}.
\]
At every old endpoint, the fact that $g_n\in C^n(\R)$ gives
\[
\alpha_j=P_J^{(j)}(u),
\qquad
\beta_j=P_J^{(j)}(v)
\qquad(0\leq j\leq n).
\]
At the new endpoints $c$ and $d$, the discrepancies in these lower-order
jets are bounded by
\[
\norm{Q_{n+1}-P_*}_{C^n([c,d])}.
\]

Apply Lemma~\ref{lem:stitch}, with $r=n$ and tolerance $\eta_n$, to each
of the finitely many intervals in $\mathcal J_{n+1}$. The corresponding
threshold depends only on the parent polynomial, the interval, and
$\eta_n$, and, crucially, not on the prescribed order-$(n+1)$ endpoint
data. Taking the minimum of these finitely many thresholds gives
$\delta_n>0$ such that all the required stitchings are possible whenever
\begin{equation}\label{eq:Qclose}
\norm{Q_{n+1}-P_*}_{C^n([c,d])}
<\min\{\delta_n,\eta_n\}.
\end{equation}
By Lemma~\ref{lem:dense-field}, choose $Q_{n+1}\in F[x]$ satisfying
\eqref{eq:Qclose}, with degree greater than $\deg P_*$ and $n+1$, and greater than the degree of every previously assigned polynomial $Q_j$.

For each $K\in\mathcal J_{n+1}$, Lemma~\ref{lem:stitch} now yields a real
polynomial $S_K$ having exactly the prescribed endpoint jets through order
$n+1$ and satisfying
\[
\norm{S_K-P_J}_{C^n(K)}<\eta_n,
\]
where $J$ is the parent of $K$. Define
\[
g_{n+1}(x)=
\begin{cases}
0, & x\notin[a,b],\\
Q_j(x), & x\in U_j,\quad 1\leq j\leq n+1,\\
S_K(x), & x\in K,\quad K\in\mathcal J_{n+1}.
\end{cases}
\]
The endpoint jets agree through order $n+1$, hence
$g_{n+1}\in C^{n+1}(\R)$. The old gap pieces are unchanged; on the new
gap, \eqref{eq:Qclose} controls $g_{n+1}-g_n$; and on each residual
interval, the same control follows from the stitching estimate. Therefore
\eqref{eq:increment} holds, and the induction is complete.



Fix $r\geq0$. For $n\geq r$, \eqref{eq:increment} together with the fact that $g_{n+1}-g_n$ vanishes outside $[a,b]$, gives
\[
\norm{g_{n+1}^{(r)}-g_n^{(r)}}_{L^\infty(\R)}=\norm{g_{n+1}^{(r)}-g_n^{(r)}}_{L^\infty([a,b])}
\leq2^{-n-2}.
\]
Thus $(g_n^{(r)})_{n\ge r}$ is uniformly Cauchy on $\mathbb R$; denote its limit by $h_r$. On every compact interval, the tails
\[
(g_n^{(r)})_{n\ge r+1}
\qquad\text{and}\qquad
(g_n^{(r+1)})_{n\ge r+1}
\]
converge uniformly to $h_r$ and $h_{r+1}$, respectively. Moreover,
\[
\bigl(g_n^{(r)}\bigr)'=g_n^{(r+1)}
\qquad (n\ge r+1).
\]
The standard theorem on uniform convergence of derivatives therefore gives
\[
h_{r+1}=h_r'.
\]
Induction on $r$ shows that $g:=h_0$ belongs to $C^\infty(\R)$ and that
$g^{(r)}=h_r$ for every $r\geq0$. Since every $g_n$ vanishes outside
$[a,b]$, the same is true of $g$, so $g\in C_c^\infty(\R)$.

Fix a gap $U_j$. The polynomial $Q_j$ assigned to $U_j$ is left unchanged from stage $j$ onward,
and therefore
\[
g|_{U_j}=Q_j|_{U_j},
\qquad
Q_j\in F[x].
\]
On the two unbounded components of $\R\setminus E$, the function $g$ is
identically zero. This proves \textup{(i)} and shows that $g$ is
real-analytic on $\R\setminus E$.

It remains to prove non-analyticity on $E$. Suppose, to the contrary,
that $g$ is real-analytic at some $x\in E$. Then $g$ is real-analytic on
some connected open interval $V\ni x$. By Lemma~\ref{lem:gaps}, there
exist two distinct gaps $U_j,U_k$ whose closures are contained in $V$.
Since $g=Q_j$ on $U_j$, the real-analytic identity theorem gives
$g=Q_j$ throughout $V$. In particular, $Q_j=Q_k$ on $U_k$, and hence
$Q_j=Q_k$ as polynomials. This contradicts
$\deg Q_j\neq\deg Q_k$. Thus $g$ is not real-analytic at any point of
$E$, proving \textup{(iii)} and completing the proof.
\end{proof}

\end{document}
